\documentclass[a4paper,11pt]{article}

% ---------------------------------------------------------------
% Packages
% ---------------------------------------------------------------
\usepackage[T1]{fontenc}
\usepackage[utf8]{inputenc}
\usepackage{lmodern}
\usepackage[margin=2.5cm]{geometry}
\usepackage{microtype}
\usepackage{booktabs}
\usepackage{tabularx}
\usepackage{multirow}
\usepackage{array}
\usepackage{graphicx}
\usepackage{xcolor}
\usepackage{listings}
\usepackage{enumitem}
\usepackage{hyperref}
\usepackage{amsmath}
\usepackage{url}
\usepackage{parskip}

\hypersetup{
  colorlinks=true,
  linkcolor=blue,
  urlcolor=blue,
  citecolor=blue
}

% ---------------------------------------------------------------
% Title and authors
% ---------------------------------------------------------------
\title{\textbf{BIU NLP at HIPE 2026:}\\
Zero/Few-Shot LLM Prompting for Person--Place\\
Relation Qualification in Historical Newspapers}

\author{
  [Author Name]\\
  Department of Computer Science\\
  Bar-Ilan University\\
  Ramat Gan, Israel\\
  \texttt{[email@biu.ac.il]}
}

\date{}

% ---------------------------------------------------------------
\begin{document}
\maketitle

% ---------------------------------------------------------------
\begin{abstract}
We describe the BIU NLP system submitted to the HIPE 2026 shared task on
person--place relation qualification in multilingual historical newspapers.
Our approach uses zero/few-shot prompting of open-weight large language models
(LLMs) without any task-specific fine-tuning. We evaluate four models ---
Gemma-4-26B (MoE), Gemma-3-27B, Mistral-Small-24B, and Aya-Expanse-32B ---
under varying prompt languages (English, native) and few-shot counts (0, 3, 5,
10). For each test file, we select the best configuration per model based on
development set macro-recall, and submit the top three models as separate runs.
Our best submission (Gemma-3-27B, English prompt, zero-shot) ranks \textbf{9th
out of 46} for French and \textbf{19th out of 47} for English in the main
evaluation. In the generalization setting (surprise French test set),
Gemma-4-26B ranks \textbf{6th out of 46} (main evaluation) and \textbf{3rd out
of 44} (binary evaluation).
\end{abstract}

% ---------------------------------------------------------------
\section{Introduction}

HIPE 2026 is a shared task at CLEF 2026 focused on \textbf{person--place
relation extraction and qualification} in digitized historical European
newspapers. Building on the entity recognition and linking tasks of HIPE-2020
and HIPE-2022, HIPE 2026 targets the question \textit{Who was where, when?} ---
aiming to support the reconstruction of life trajectories and mobility patterns
from historical sources.

The task presents several challenges that make standard NLP approaches
difficult to apply directly: (1)~text is in three languages (German, English,
French) spanning the 18th to early 20th century; (2)~documents are digitized
from microfilm and contain significant OCR noise; (3)~the relations are
semantically nuanced and context-dependent, not reducible to simple entity
co-occurrence; and (4)~training data is limited relative to the number and
diversity of newspapers covered.

Given these characteristics, we hypothesized that large language models with
strong multilingual and reasoning capabilities, prompted with explicit task
definitions and label constraints, could compete effectively without the
overhead of fine-tuning. This paper describes our system, experimental
results, and analysis.

% ---------------------------------------------------------------
\section{Related Work}

\paragraph{Historical Document NLP.}
Processing historical documents poses unique challenges compared to
contemporary text, most notably OCR noise, archaic spelling, and the absence
of large annotated corpora. Prior work has addressed these through
normalization pipelines, domain-adaptive pre-training, and data augmentation.
The HIPE shared tasks (2020, 2022) established benchmarks for named entity
recognition and entity linking in historical newspapers across multiple
European languages, with participating systems relying primarily on fine-tuned
transformer models such as BERT and its multilingual and domain-specific
variants.

\paragraph{Relation Extraction.}
Relation extraction (RE) has a long history in NLP, from early pattern-based
and kernel methods to neural sequence labeling and graph-based approaches.
Most modern RE systems are trained on annotated datasets such as ACE, TACRED,
or DocRED. However, these resources cover contemporary English and do not
transfer well to historical multilingual text with noisy OCR. Supervised
cross-lingual RE has been explored using multilingual pre-trained models, but
requires parallel or translated training data that is often unavailable for
historical domains.

\paragraph{LLM-based Information Extraction.}
Recent work has demonstrated that large language models can perform
competitive zero-shot and few-shot information extraction without fine-tuning.
In-context learning with structured output formats (e.g., JSON) has been shown
to be effective for named entity recognition, relation extraction, and event
detection. Chain-of-thought prompting, which elicits explicit intermediate
reasoning before a final answer, has been shown to improve performance on
tasks requiring multi-step inference, and is particularly relevant for
nuanced relation classification where evidence must be grounded in specific
text spans.

\paragraph{Multilingual LLMs.}
The models we evaluate span different strategies for multilingual competence.
Gemma-3 and Gemma-4 are general-purpose models with strong multilingual
capabilities due to their large pre-training corpora. Mistral Small was
optimized for efficiency across European languages. Aya-Expanse is explicitly
designed as a multilingual instruction-following model, covering over 23
languages including all three task languages, and has shown strong performance
on cross-lingual benchmarks. Mixture-of-Experts (MoE) architectures, as used
in Gemma-4, activate only a fraction of parameters per token, enabling large
total parameter counts at lower inference cost --- making them attractive for
long-context historical document processing.

\paragraph{Positioning.}
Our work contributes to the emerging line of research on LLM-based historical
RE by (1)~comparing multiple open-weight models under a controlled experimental
grid, (2)~investigating the effect of prompt language (English vs.\ native) in
a multilingual historical setting, and (3)~analyzing the discrepancy between
in-domain and out-of-domain generalization for MoE vs.\ dense models.

% ---------------------------------------------------------------
\section{Task Description}

The HIPE 2026 relation qualification task takes as input a historical newspaper
article with pre-annotated named entities (persons and locations linked to
Wikidata where possible) and a set of candidate \textbf{person--place pairs}
extracted from the article. For each pair, the system must assign labels for
two relations:

\begin{itemize}[noitemsep]
  \item \textbf{\texttt{at}}: Does the article provide textual evidence that
        the person was at the location at any point before publication?\\
        Labels: \texttt{TRUE}, \texttt{PROBABLE}, \texttt{FALSE}.
  \item \textbf{\texttt{isAt}}: Does the article support that the person's
        location in the article's immediate temporal horizon (current, ongoing,
        or very recent) is the given place?\\
        Labels: \texttt{TRUE}, \texttt{FALSE}.
\end{itemize}

The two relations are constrained: if \texttt{at} is \texttt{FALSE} then
\texttt{isAt} must also be \texttt{FALSE}; if \texttt{isAt} is \texttt{TRUE}
then \texttt{at} must also be \texttt{TRUE}. The label \texttt{PROBABLE} is
only valid for \texttt{at}.

The evaluation metric is \textbf{global macro-recall}, averaged equally over
both relations. A system predicting all pairs as \texttt{FALSE} achieves
approximately \textbf{0.4167} macro-recall (the class-imbalance baseline),
since \texttt{FALSE} is the dominant label for both relations. A secondary
\textbf{binary evaluation} maps \texttt{PROBABLE} predictions for \texttt{at}
to \texttt{TRUE}, simplifying the problem to binary classification.

% ---------------------------------------------------------------
\section{Data}

The HIPE 2026 dataset is derived from the Impresso project, which digitized
and annotated historical European newspapers. The data builds on the HIPE-2022
v2.1 named entity-annotated corpus, covering multiple newspaper collections in
German, English, and French spanning roughly 1780--1950. Named entity
annotations include Wikidata entity identifiers for persons and locations.
Person--place pairs were extracted from each article, pre-annotated by an LLM
ensemble, and then manually reviewed and corrected.

\subsection{Dataset Statistics}

Table~\ref{tab:data} shows the number of articles and annotated person--place
pairs per language and split used in our experiments.

\begin{table}[h]
\centering
\caption{Dataset statistics (articles~/~annotated pairs).}
\label{tab:data}
\begin{tabular}{lrrrr}
\toprule
\textbf{Split} & \textbf{German} & \textbf{English} & \textbf{French} & \textbf{Total} \\
\midrule
Train          & 88~/~1{,}224   & 56~/~496         & 317~/~4{,}450   & 461~/~6{,}170  \\
Dev            & 32~/~432       & 17~/~151         & 107~/~1{,}498   & 156~/~2{,}081  \\
Test (impresso)& 19~/~238       & 19~/~162         & 19~/~238        & 57~/~638       \\
Test (surprise-FR) & ---        & ---              & 30~/~480        & 30~/~480       \\
\bottomrule
\end{tabular}
\end{table}

The \textbf{surprise test set} is drawn from a different corpus (French
literary works) not revealed before submission, designed to test
out-of-domain generalization. The label distribution is heavily skewed toward
\texttt{FALSE} for both relations, reflecting that most person--place pairs in
a given article have no meaningful spatial relationship.

\subsection{Data Characteristics}

Key challenges include: (1)~OCR noise from historical digitization, including
hyphenation artifacts, spelling variants, and line-break insertions within
tokens; (2)~archaic and historical spelling, particularly in German and older
French; (3)~long documents containing many persons and places, requiring the
model to scope its attention to the specific pair; and (4)~ambiguous entity
mentions appearing in different surface forms across the same article.

% ---------------------------------------------------------------
\section{Approach}

\subsection{Overview}

We frame the task as a \textbf{per-pair text classification} problem addressed
via zero/few-shot prompting of large language models, with no task-specific
fine-tuning. Each person--place pair is classified independently: the model
receives the full article text and the two entity mention strings, and outputs
a structured prediction.

\subsection{Prompt Design}

Prompts instruct the model to act as a historian reading a historical European
newspaper article. Each prompt contains six components:

\begin{enumerate}[noitemsep]
  \item \textbf{System role}: ``You are a historian working on person-location
        relations in multilingual historical European newspapers.''
  \item \textbf{Relation definitions}: explicit temporal scope and label
        constraints for \texttt{at} and \texttt{isAt}.
  \item \textbf{Decision guidance}: rules specifying when each label applies
        (e.g., \texttt{PROBABLE} only for indirect or partial evidence;
        \texttt{isAt} is never \texttt{PROBABLE}; logical constraints between
        the two relations).
  \item \textbf{Robustness rules}: explicit instructions to handle OCR noise,
        line-break artifacts, and historical spelling without external knowledge.
  \item \textbf{Chain-of-thought format}: the model first outputs a brief
        explanation for each relation (\texttt{at\_explanation},
        \texttt{isAt\_explanation}, at most 100 words each, citing only
        article evidence) before giving the label.
  \item \textbf{Structured JSON output}: the final answer is a JSON object
        with keys \texttt{at\_explanation}, \texttt{at},
        \texttt{isAt\_explanation}, \texttt{isAt}.
\end{enumerate}

The chain-of-thought explanation encourages the model to ground decisions in
the article text and provides interpretable traces for error analysis.
Prompts were prepared in \textbf{English} and in the \textbf{native document
language} (German, French). We hypothesized that native-language prompts might
benefit models less robustly trained on historical German.

\subsection{Few-Shot Examples}

We experimented with 0, 3, 5, and 10 in-context examples per prompt.
Examples were sampled from the training set to cover a mix of \texttt{TRUE},
\texttt{PROBABLE}, and \texttt{FALSE} labels, avoiding label bias. Each
example includes the full JSON output with explanations. At 10-shot, prompts
reach several thousand tokens; we verified that all models supported the
required context lengths.

\subsection{Models}

We evaluated four open-weight instruction-tuned models, all loaded via
HuggingFace Transformers in bfloat16 precision on NVIDIA A100-SXM4-80GB GPUs.
Table~\ref{tab:models} summarizes the models.

\begin{table}[h]
\centering
\caption{Models evaluated.}
\label{tab:models}
\begin{tabular}{llrl}
\toprule
\textbf{Model key} & \textbf{HuggingFace ID} & \textbf{Params} & \textbf{Architecture} \\
\midrule
gemma4   & google/gemma-4-26B-A4B-it                  & 26B (4B active) & MoE   \\
gemma3   & google/gemma-3-27b-it                      & 27B             & Dense \\
mistral  & mistralai/Mistral-Small-24B-Instruct-2501  & 24B             & Dense \\
aya      & CohereForAI/aya-expanse-32b                & 32B             & Dense, multilingual \\
\bottomrule
\end{tabular}
\end{table}

Gemma-4 is a Mixture-of-Experts model activating only 4B parameters per
forward pass despite its 26B total, making it computationally efficient.
Aya-Expanse is purpose-built for multilingual tasks and natively covers all
three task languages.

\subsection{Configuration Selection and Submission Strategy}

For each model $\times$ prompt language $\times$ shot count combination, we
ran inference on the development set and computed global macro-recall. We
selected the best configuration per model per language and submitted the top-3
models as run1/run2/run3, yielding 12 submission files (3~runs $\times$
4~test files). Table~\ref{tab:submissions} shows the final configurations.

\begin{table}[h]
\centering
\caption{Submitted configurations per test file.}
\label{tab:submissions}
\setlength{\tabcolsep}{5pt}
\begin{tabular}{llll}
\toprule
\textbf{Test file} & \textbf{run1} & \textbf{run2} & \textbf{run3} \\
\midrule
impresso-test-de   & gemma4 / en / 3-shot   & mistral / native / 0-shot & gemma3 / native / 0-shot \\
impresso-test-en   & gemma4 / en / 5-shot   & gemma3 / en / 0-shot      & mistral / native / 0-shot \\
impresso-test-fr   & gemma4 / en / 3-shot   & gemma3 / en / 0-shot      & mistral / native / 0-shot \\
surprise-test-fr   & gemma4 / en / 3-shot   & gemma3 / en / 0-shot      & mistral / native / 0-shot \\
\bottomrule
\end{tabular}
\end{table}

% ---------------------------------------------------------------
\section{Results}

\subsection{Development Set Results}

Table~\ref{tab:dev} reports the best development set macro-recall per model
across all configurations. Gemma-4 consistently achieves the highest scores,
with 10-shot English prompts performing best across all three languages.

\begin{table}[h]
\centering
\caption{Best development set results (global macro-recall). All-FALSE baseline: $\approx$0.4167.}
\label{tab:dev}
\begin{tabular}{lrrrp{4cm}}
\toprule
\textbf{Model} & \textbf{DE} & \textbf{EN} & \textbf{FR} & \textbf{Best config} \\
\midrule
gemma-4-26B (MoE)  & \textbf{0.7518} & \textbf{0.7803} & \textbf{0.7481} & en / 10-shot \\
gemma-3-27b        & 0.5622 & 0.7209 & 0.6611 & native/0-shot (DE); en/0-shot (EN, FR) \\
Mistral-Small-24B  & 0.5917 & 0.5935 & 0.5914 & native / 0-shot \\
aya-expanse-32b    & 0.5950 & 0.5967 & 0.6183 & en / 0-shot \\
\bottomrule
\end{tabular}
\end{table}

All models substantially exceed the all-FALSE baseline. Gemma-4 shows a
particularly large margin for English (+0.36) and French (+0.33). Native-
language prompting generally benefits Mistral on German and Gemma-3 slightly
on German, while English prompts perform best for French and English across
all models.

\subsection{Effect of Shot Count}

Increasing in-context examples generally improved Gemma-4 up to 10~shots
(e.g., $+$0.12 macro-recall for DE from 0-shot to 10-shot). For Gemma-3 and
Mistral, few-shot examples were inconsistent and sometimes degraded
performance, possibly due to prompt length effects or reduced attention to the
test article. We therefore submitted zero-shot configurations for these models.

\subsection{Official Test Results}

Tables~\ref{tab:test-main} and~\ref{tab:test-binary} report official test
results from the HIPE 2026 evaluation. Rankings are shown as rank~/~total
submissions (46--47 impresso, 44 surprise submissions).

\begin{table}[h]
\centering
\caption{Main evaluation --- impresso profile score (macro-recall). Format: rank/total --- score.}
\label{tab:test-main}
\setlength{\tabcolsep}{4pt}
\begin{tabular}{lrrr}
\toprule
\textbf{Profile} & \textbf{run1 (gemma4)} & \textbf{run2 (gemma3)} & \textbf{run3 (mistral)} \\
\midrule
Overall (mean DE/EN/FR) & 41/46 --- 0.4429 & \textbf{20/46 --- 0.5781} & 24/46 --- 0.5390 \\
German                  & 37/46 --- 0.4495 & 31/46 --- 0.5050          & \textbf{20/46 --- 0.5451} \\
English                 & 36/47 --- 0.4583 & \textbf{19/47 --- 0.5762} & 28/47 --- 0.5101 \\
French                  & 42/46 --- 0.4208 & \textbf{9/46 --- 0.6529}  & 23/46 --- 0.5617 \\
Surprise-FR             & \textbf{6/46 --- 0.6837}  & 22/46 --- 0.5265 & 17/46 --- 0.6000 \\
\bottomrule
\end{tabular}
\end{table}

\begin{table}[h]
\centering
\caption{Binary evaluation (\texttt{PROBABLE} ``at'' $\to$ \texttt{TRUE}). Format: rank/total --- score.}
\label{tab:test-binary}
\setlength{\tabcolsep}{4pt}
\begin{tabular}{lrrr}
\toprule
\textbf{Profile} & \textbf{run1 (gemma4)} & \textbf{run2 (gemma3)} & \textbf{run3 (mistral)} \\
\midrule
Overall (mean DE/EN/FR) & 41/46 --- 0.5252 & \textbf{20/46 --- 0.6721} & 24/46 --- 0.6274 \\
German                  & 40/46 --- 0.5329 & 24/46 --- 0.6231          & \textbf{23/46 --- 0.6381} \\
English                 & 36/47 --- 0.5382 & \textbf{21/47 --- 0.6524} & 28/47 --- 0.5904 \\
French                  & 42/46 --- 0.5044 & \textbf{14/46 --- 0.7409} & 23/46 --- 0.6538 \\
Surprise-FR             & \textbf{3/44 --- 0.8804}  & 26/44 --- 0.6473 & 20/44 --- 0.7262 \\
\bottomrule
\end{tabular}
\end{table}

The per-relation breakdown for our best run (run2, French impresso) is:
\texttt{at} macro-recall 0.5703 (binary: 0.7462), \texttt{isAt} macro-recall
0.7356, \texttt{at} accuracy 0.6387 (binary: 0.7227), \texttt{isAt} accuracy
0.7773.

% ---------------------------------------------------------------
\section{Analysis and Discussion}

\subsection{Gemma-3 as the Most Reliable System}

Gemma-3 (run2) is our strongest overall submission, ranking 20/46 and
achieving the best individual result among our runs (French, rank~9/46). Its
consistent zero-shot English performance across all three languages makes it a
strong approach for multilingual historical RE. The strong French result may
partly reflect the larger French training set (317 articles, 4,450 pairs),
which provides richer in-domain calibration data for selecting examples and
verifying predictions at development time.

\subsection{The Gemma-4 Anomaly}

The most striking finding is the divergence between Gemma-4's development and
test performance. On the dev set it was the clear winner (0.75--0.78
macro-recall), yet on the impresso test sets it ranked last among our runs
(41/46 overall). Conversely, on the surprise French test it ranked 6/46 main
and 3/44 binary --- best among all three of our submissions for that set.

Several hypotheses may explain this divergence:

\begin{enumerate}[noitemsep]
  \item \textbf{Dev/test domain shift within impresso}: the dev and test splits
        may differ in newspaper titles or time periods, and Gemma-4 may have
        overfit to dev-set idiosyncrasies rather than generalizable patterns.
  \item \textbf{\texttt{PROBABLE} label calibration}: Gemma-4 may predict
        \texttt{PROBABLE} more frequently, which benefits macro-recall on the
        dev set but backfires if the test distribution is less ambiguous. The
        binary evaluation, which maps \texttt{PROBABLE}~$\to$~\texttt{TRUE},
        shows a milder drop for Gemma-4 (0.52 vs.\ 0.44 in main), consistent
        with this.
  \item \textbf{Few-shot example bias}: the 3/5-shot examples used for impresso
        test submissions may have introduced distribution shifts absent from
        the dev set.
  \item \textbf{Surprise set characteristics}: the literary works corpus may be
        stylistically closer to Gemma-4's training distribution, explaining
        its strong out-of-domain generalization.
\end{enumerate}

\subsection{Prompt Language Effects}

Native-language prompting helped Mistral most on German (0.5917 native
vs.\ 0.5389 en, 0-shot), consistent with Mistral Small's training being less
robust on historical German text. For French, native prompts did not help any
model meaningfully, likely because English and French are both well-represented
in all model pretraining corpora. For Gemma-3, native prompts provided a
modest benefit on German (0.5622 native vs.\ 0.5299 en).

\subsection{Few-Shot Scaling}

Few-shot examples improved Gemma-4 substantially but were inconsistent for
Gemma-3 and Mistral. This may reflect differences in instruction-following
capacity: Gemma-4's larger capacity allows it to integrate more examples
effectively, while smaller models may lose focus on the target article as
prompt length grows. Retrieval-augmented example selection (e.g., choosing
examples similar in article topic or entity type to the target pair) could
address this.

\subsection{Progress Beyond Baseline}

All submitted runs exceed the all-FALSE baseline ($\approx$0.4167). Run2
achieves 0.5781 overall (main) and 0.6721 (binary), gains of $+$0.16 and
$+$0.26 respectively. The French result (0.6529, rank~9/46) represents a gain
of $+$0.24, and the surprise-FR binary result (0.8804, rank~3/44) a gain of
$+$0.46. These results demonstrate that zero/few-shot LLM prompting is a
strong approach for this task even without any fine-tuning.

% ---------------------------------------------------------------
\section{Perspectives for Future Work}

\textbf{Fine-tuned smaller models.}
The strong zero-shot performance of 24--27B models motivates exploring whether
a smaller model (7B--13B) fine-tuned on the HIPE 2026 training set could match
this performance at lower inference cost.

\textbf{Retrieval-augmented few-shot.}
Rather than fixed in-context examples, selecting examples by similarity to the
target pair (article topic, entity type, time period) could improve few-shot
performance across all models.

\textbf{Constrained decoding.}
The logical constraints between \texttt{at} and \texttt{isAt} could be
enforced via constrained decoding or post-processing, rather than relying on
the model to satisfy them through prompting alone.

\textbf{Cross-lingual transfer.}
The French training set is roughly $6\times$ larger than German and $9\times$
larger than English. Systematic cross-lingual transfer experiments (training
on French, evaluating on German/English) could address the data imbalance and
improve low-resource language performance.

\textbf{Analysis of the Gemma-4 gap.}
A detailed comparison of Gemma-4 predictions on dev vs.\ test --- examining
label distributions, entity types, and article domains --- would clarify
whether the dev/test discrepancy is driven by calibration, domain mismatch,
or few-shot example selection.

% ---------------------------------------------------------------
\section*{Acknowledgments}

This work was supported by the Israeli Innovation Authority.

% ---------------------------------------------------------------
\section*{Declaration on Generative AI}

Claude Code was used to assist with coding and implementation. Gemini was used for rephrasing and editorial assistance.

% ---------------------------------------------------------------
% References will be added here
% \bibliographystyle{plain}
% \bibliography{references}

\end{document}